\chapter{最后那段有限计算，到底在调整什么}\label{sec:controllers}
\section{只知道比例正确，还差一步}
模五的 $(1,1)$ 与 $(2,2)$ 比例相同，但列并不相同。上一讲的精确输送要求\emph{整列}同余。因此，先把比例调对以后，还得说明允许怎样改变共同倍数。

这不是要求实际数值变大。例如某段平均可以输出一个很小的共同倍数；把这个倍数从分析代表中取出后，本原列却可能变成另一个有限尺度。我们始终跟踪实际完整配置和分析代表之间的对应。

\section{一张整数矩阵，可以产生几种本原输出}
设输入列 $(a,b)$ 的最大公因数为一。施加整数矩阵后，两坐标可能有公因子。把这个最大公因数叫作\textbf{内容}，记为 $\cont$。

最简单的例子是 $A=\diag(1,6)$。输入 $(s,t)$ 本原，输出 $(s,6t)$ 的公因子是 $\gcd(s,6)$，所以只可能为 $1,2,3,6$。

甚至可以固定输入模五都等于 $(1,1)$，而让四种内容分别出现：
\[
\begin{array}{c|c|c}
\text{输入}&A\text{ 的输出}&\text{本原输出}\\ \hline
(1,1)&(1,6)&(1,6)\\
(16,1)&(16,6)&(8,3)\\
(21,1)&(21,6)&(7,2)\\
(6,1)&(6,6)&(1,1)
\end{array}
\]
它们分别除去了 $1,2,3,6$。所以同一个模五输入类，在施加同一个 $A$ 并取本原代表以后，会得到几种不同列。这就是“内容分支”的具体意思。
\begin{example}[一个内容分支的具体样子]
令 $A=\diag(1,6)$，控制模数取 $q=5$。
对本原列 $(s,t)^T$，内容为 $\gcd(s,6)$，可能取 $1,2,3,6$。
固定任意本原模 $5$ 类，在模 $2,3$ 处分别选择 $s$ 是否为零，并取 $t$ 为单位，
就能安排这四种内容；与模 $5$ 的条件由中国剩余定理合并。
本原输出模 $5$ 因而分别是 $Ax,2^{-1}Ax,3^{-1}Ax,6^{-1}Ax$。
一般矩阵先经 Smith 标准形化为这个例子的形式。
\end{example}
\begin{lemma}[本原化的全部有限分支]\label{lem:content}
设 $A\in M_2(\Z)$ 满秩，条目最大公因子为 $1$，$d=|\det A|$、$\gcd(d,q)=1$。固定模 $q$ 本原列 $x$，则所有整数本原提升的本原输出模 $q$ 恰为
\begin{equation}\label{eq:content}
\{g^{-1}Ax:g\mid d,\ g>0\}.
\end{equation}
若 $\Gamma(q,\Z)$ 已实现，每个分支都能从当前列实现。
\end{lemma}
\begin{proof}
伴随恒等式说明 $\cont(Av)\mid d$。写 Smith 分解 $A=U\diag(1,d)V$；对本原 $Vv=(s,t)$，内容为 $\gcd(s,d)$。给定 $g\mid d$，在每个 $\ell^e\Vert d$ 指定 $s\equiv\ell^{v_\ell(g)}\pmod{\ell^e}$（指数为 $e$ 时取 $0$）、$t\equiv1\pmod{\ell^e}$。与模 $q$ 数据用 CRT 合并。所得本原剩余列可提升为本原整数列：固定非零第二坐标后，在第一坐标的等差类中避开第二坐标所含的其余有限素因子。其内容恰为 $g$。最后用引理\ref{lem:fiber}实现准备。
\end{proof}
该引理为什么要先有 $\Gamma(q,\Z)$？因为我们需要先把当前实际列送到具有指定内容的那个提升，再执行 $A$。没有这个先行准备，就不能把“某个提升存在”当成“当前输入能走过去”。

\section{有限集合上的正向操作为什么可以组成群}
假设一个置换在六步后回到原样，那么它的逆就是再做五步。这是有限群中最朴素的事实。在本题中，每条安全内容分支是有限标记集上的置换，所以无需额外要求它的有理矩阵具有统一逆词，仍能在这个有限集合里反向移动。

更妙的是，先走一次内容为 $g$ 的分支，再走若干次内容一分支，将矩阵本身绕回去，剩下的作用就是乘 $g^{-1}$。因此行列式的素因子可以给有限尺度变化。
\begin{corollary}[有限正向对应的群化]\label{cor:finite-group}
设若干返回 $A_i$ 满足引理\ref{lem:content}，并在模 $q$ 保持一个对单位标量封闭的合法列集。它们的正向内容对应轨道等于
\[
\langle A_i\bmod q,\ \ell I:\ell\mid\prod_i\det A_i\text{ 为素数}\rangle
\]
的群轨道。
\end{corollary}
\begin{proof}
每个分支 $g^{-1}A_i$ 是有限集上的置换。内容 $1$ 分支给 $A_i$，其逆为某个非负幂。若 $A_i$ 的阶为 $m$，一次内容 $g$ 分支后接 $m-1$ 次内容 $1$ 分支，给标量 $g^{-1}I$。这些恰生成所示群。
\end{proof}

\section{从模三走到模九：为什么有时会卡住}
普通平移 $x\mapsto x+1$ 在模 $\ell^b$ 下当然走遍所有类。一般分式变换模 $\ell$ 看起来像平移，却可能在更高精度过早循环。因此要额外检查每走 $\ell$ 步，步长是不是恰好多出一个 $\ell$ 因子。
\begin{definition}[局部整数]
$\Z_{(\ell)}$ 是分母不被 $\ell$ 整除的有理数环。
它的单位是赋值为零的元素，可同时约化到所有 $\Z/\ell^b\Z$。
当列的第二坐标是单位时，用其比值 $x=z/y$ 表示仿射方向。
矩阵 $\left(\begin{smallmatrix}a&b\\c&d\end{smallmatrix}\right)$ 对方向的作用为
$x\mapsto(ax+b)/(cx+d)$，只要分母为单位。
\end{definition}
若 $N^2=cI$，所有高次幂都能写成 $A I+B N$，因为偶次幂给标量，奇次幂给 $N$。于是很多矩阵计算只剩两项。例如
\[(I+N)^3=(1+3c)I+(3+c)N.
\]
当 $c=0\pmod3$，$3+c$ 通常恰含一个三；但 $c=6\pmod9$ 时它含至少两个三，就可能跳过一层。这解释了后面为什么会特意避开一个模九余类。
\begin{lemma}[局部平移的逐层提升]\label{lem:local-cycle}
设 $T$ 在 $\Z_{(\ell)}$ 上整可逆，模 $\ell$ 的射影作用为非平凡平移。写 $T=a(I+N)$、$N^2=cI$。若 $\ell\ge5$，或 $\ell=3$ 且 $c\not\equiv6\pmod9$，则其在每个 $\Z/\ell^b\Z$ 的仿射方向上为单循环。对 $\ell=2$，同一结论在 $a$ 为奇单位、$c$ 为偶数时成立。
\end{lemma}
\begin{proof}
此时 $c\in\ell\Z_{(\ell)}$。展开 $(I+N)^\ell=A I+B N$；$A$ 为单位，$v_\ell(B)=1$。在 $\ell=3$ 时 $B=3+c$，在 $\ell=2$ 时 $A=1+c,B=2$。反复取 $\ell$ 次幂后，规范化的 $N$ 系数赋值每次增加一。一般幂的系数赋值因而等于指数的 $\ell$ 进赋值。其在任一方向 $x$ 上的位移还乘一个模 $\ell$ 等于原平移步长的单位。因此固定 $x\bmod\ell^b$ 的指数恰为 $\ell^b$ 的倍数，轨道长度为 $\ell^b$。
\end{proof}

\section{我们到底要把参数送到哪里}
回到 $K_{p,r}$，令 $y=a+pz$，合法性为 $\gcd(y,n)=1$。想选 $j+1$ 份第一块、$j$ 份第二块，以及 $p-1-2j$ 份载体。将它们的和算出来：
\[
(j+1)a+j(rz-a)-(p-1-2j)pz=y-(p^2-nj)z.
\]
所以只要把参数方向送到
\begin{equation}\label{eq:terminal}
t_j=\binom{1}{p^2-nj},\qquad
\max(0,\lceil(p-1-r)/2\rceil)\le j\le(p-1)/2,
\end{equation}
这组数的和就是零。$j$ 的区间只是要求载体取用数在 $0,r$ 之间。

举例：$p=5,n=11,r=1$，区间只允许 $j=2$，终端列是 $(z,y)=(1,3)$。于是 $a=y-pz=-2$，第二块值为 $3$，载体为 $-5$。三份 $-2$ 加两份 $3$ 恰好为零。这是一个非零的配置，但已经能接第5讲的收尾。

\section{为什么只会生成少量小数字，也够用了}
互补混合中，份数最多是 $p$。因此容易得到的是不超过 $p$ 的安全数字，而不是模 $n$ 的每个单位直接给一个操作。下一个引理说，在 $2p<n<4p$ 这段范围，所有小安全数字已经能通过相乘和取逆生成模 $n$ 的全部单位。

证明不是找原根。奇数 $n$ 时，假设有一个最小缺失单位，通过乘二或四再折回较小正代表，会得到更小的缺失单位，矛盾。偶数时则按是否被四整除，利用平衡代表或数量超过群的一半来完成。
\begin{lemma}[小数字生成模n单位群]\label{lem:small-units}
设 $p\ge5$、$2p<n<4p$、$\gcd(p,n)=1$。$-1$ 和所有 $1\le b\le p$、$\gcd(b,n)=1$ 的类生成 $(\Z/n\Z)^\times$。若 $n$ 偶，还存在安全奇数 $b\le p$ 满足 $b\equiv3$ 或 $5\pmod8$。
\end{lemma}
\begin{proof}
记第一项的生成群为 $H$。$n$ 奇时，若 $H$ 不满，取最小正缺失单位 $p<a<n/2$。$a<n/3$ 时取 $4a-n$，$a>n/3$ 时取 $n-2a$，均为更小的正缺失单位，矛盾。
若 $n=2m$、$m$ 奇，选奇数 $b=(m\pm1)/2\le p$，则 $H$ 模 $m$ 包含 $2$ 及其逆。每个模 $m$ 单位的平衡代表绝对值小于 $p$，除尽 $2$ 因子后即是已知奇单位。
若 $4\mid n$，单位在 $(0,n/2)$ 中由 $a\mapsto n/2-a$ 配对，所以 $(0,n/4)$ 恰有 $\varphi(n)/4$ 个单位。它们及负数已属于 $H$，再加 $\pm p$，得到 $|H|>\varphi(n)/2$。

第二项先检查 $3,5,11,13$。$p=5,7$ 时若 $3,5$ 都不安全，则 $30\mid n<4p$，矛盾；$p=11$ 时同理由 $165\mid n$ 排除。若 $p\ge13$ 且四者均不安全，记 $n$ 的不同奇素因子数为 $d\ge4$，$\rho=\prod_{\ell\mid n,\ell\text{奇}}(1-1/\ell)$。等差列 $3+8t\le p$ 的项数 $L\ge(p-4)/8$，容斥给安全项数至少 $L\rho-2^d$。而
\[
L\rho>n\rho/32-1/2\ge60\cdot6^{d-4}-1/2>2^d.
\]
这里 $n\rho\ge1920\cdot6^{d-4}$，由四个指定素因子及其余奇素因子得到。
\end{proof}

\section{把所有工具接成一条有限路径}
下面定理列出三个输入条件，你可以把它看成后面构造的任务单：已经能在同余类内部精确移动；能在模 $n$ 下改变第一坐标；还要能处理有限共同尺度和二进信息。
\begin{theorem}[三块核心的共同终端定理]\label{thm:terminal}
在\eqref{eq:core}上，令 $h=\delta n$、$q=h^2$，其中 $\delta=1$（$n$ 奇）或 $\delta\in\{2,4\}$（$n$ 偶）。假设：
\begin{enumerate}
\item $L(h\Z)$ 和 $\Gamma(q,\Z)$ 可实现；
\item 若干条目本原的整数满秩返回，其行列式与 $q$ 互素，模 $q$ 保持合法域，且有限内容群包含 $V\equiv U(-p^{-2})\pmod n$；
\item $n$ 奇时该群包含标量 $2I$；$n$ 偶时\eqref{eq:terminal}的 $j$ 区间含 $\delta$ 个连续整数。
\end{enumerate}
则全部合法核心可归零。$p=3,n=7$ 时同样成立。
\end{theorem}
\begin{proof}

\par\noindent\textbf{第一步：获得模 $n$ 的全部有限标量。}
保留载体，两个 $p$ 块互补混合给 $(a,z)$ 坐标中的
\[
p^{-1}P_b,\qquad P_b=\begin{pmatrix}b&r(p-b)/2\\0&p\end{pmatrix},
\quad b\text{ 奇},\quad1\le b\le p.
\]
当 $\gcd(b,n)=1$ 时，$P_b$ 的行列式与 $q$ 互素，且 $y'\equiv by\pmod n$。$b<p$ 时条目本原，故推论\ref{cor:finite-group}提供 $p,b$ 的素因子标量；$P_1$ 已提供 $p$。结合假设中的 $2I$ 及引理\ref{lem:small-units}，有限标量模 $n$ 满射。

\par\noindent\textbf{第二步：先把列送对模 $n$。}
因为 $y$ 为单位，$z/y$ 有定义。
矩阵 $V\equiv U(-p^{-2})$ 把比值减去 $p^{-2}$，故可以一次指定所需平移次数。
对当前列，取 $k\equiv p^2zy^{-1}-1\pmod n$，用 $V^k$ 后再乘一个可用有限标量，得到 $(z,y)\equiv(1,p^2)\pmod n$。


\par\noindent\textbf{第三步：处理额外的二进信息。}
$n$ 奇时，$h=n$，本步骤无需操作，直接进入下一步。
$n$ 偶时写 $n=2^a m$、$m$ 奇，令 $B=2^{2v_2(h)}$。有限群中的 $P_bP_1^{-1}$ 及标量给
\begin{equation}\label{eq:split}
D_b=\begin{pmatrix}b^{-1}&0\\n(b^{-1}-1)/2&1\end{pmatrix}.
\end{equation}
引理\ref{lem:small-units}选出 $b\equiv3,5\pmod8$，因 $v_2(b^2-1)=3$，$-1,b$ 生成每个有限二进单位群。用\eqref{eq:split}及 $D_{-1}$ 使 $z=1\pmod B$，同时保持 $y\pmod n$。$V\bmod2$ 为上幺幂，逐次平方给 $V^B\equiv I\pmod B$；使用 $V$ 的 $B$ 倍次数恢复模 $m$ 的首坐标。于是 $z=1\pmod h$。

\par\noindent\textbf{第四步：从浅层修正到平方层。}
按 $y$ 选 $j\bmod\delta$，使当前列与 $t_j$ 模 $h$ 相同。置 $E=E_{21}$、$X=VEV^{-1}$，则模 $n$
\[
Et_j=(0,1)^T,\qquad Xt_j=(-2/p^2,2)^T.
\]

两列的行列式 $2/p^2$ 是模 $m$ 单位。
为什么这里仅需求解\emph{线性}同余？因为 $q=h^2$，对于任意矩阵 $A,B$，
\[
(I+hA)(I+hB)\equiv I+h(A+B)\pmod{h^2}.
\]
二次交叉项消失。两个独立切向列于是能修正任意两坐标的一阶误差。
在奇数部分，$2$ 可逆，所以这两个修正方向足够。用 $L(hk)$ 及其有限 $V$ 共轭，解一次二维线性同余，将两坐标修正到 $t_j\pmod{m^2}$；参数用 CRT 保持模 $B$。随后一次 $L(hk)$ 修正第二坐标模 $B$，并取 $k\equiv0\pmod m$。最终列与 $t_j$ 模 $q$ 相同。

\par\noindent\textbf{第五步：由同余相同到实际相同。}
每条有限内容分支由引理\ref{lem:content}提升，最后用引理\ref{lem:fiber}精确送到 $t_j$。平均其零和组，再用引理\ref{lem:zero-tail}收尾。$p=3,n=7$ 时，模 $7$ 的单位由 $-1,2$ 生成，其余步骤相同，收尾用引理\ref{lem:zero-base}。
\end{proof}

\section{再慢看一次：模n正确以后，误差怎样被消掉}
最容易被一句“线性修正”带过的是第四步。假设当前列 $w$ 与目标 $t$ 已模 $h$ 相同，写
\[w=t+he.
\]
施加 $I+hX$，模 $h^2$ 下得到
\[(I+hX)w\equiv t+h(e+Xt)\pmod{h^2}.
\]
于是修正整列，只需让 $Xt\equiv-e\pmod h$。如果我们手里的两个矩阵方向作用于 $t$ 后成为两条独立列，这就是一个二元线性方程组。

本题中的两条列为 $(0,1)^T$ 与 $(-2/p^2,2)^T$，行列式是 $2/p^2$。对奇素数部分可逆，所以够用；对二进部分不够，因此才先用乘法调整第一坐标，再用一次下剪切调整第二坐标。这就是定理分开处理奇偶部分的真正原因。

最后列已经与 $t_j$ 模 $h^2$ 完全相同，上一讲的精确纤维引理才出场。这时无限数值大小被处理掉，接下来平均零和组并收尾，全程只有有限次操作。
